% Pursuing Stacks (cotes 134-2 to 134-8): archive cote 134-2, site page 34 (PDF page 35).
% letter to Quillen. Draft compared: ps1.tex:990-1057.
% AI-transcribed by Claude Opus 5.5 (claude-opus-5-5, adaptive thinking, effort high), Message Batches API, batch ps1 (msgbatch_01Rv5hCaNTtUt4JHx1yqoHPL), 2026-11-12; not yet checked by eye.
\scanpage{134-2/34}
\mspage{\hand{11}}

\mspage{10} 22.2.1983

\over{\unsure{12}}{13}. It seems I can't help pursuing further the reflection I started with
this letter ! First I would like to come back upon terminology. Maybe
to give the name of n-groupoids and $\infty$-groupoids to the objects I was af-%
ter is not proper, for two reasons: a) it conflicts with a standing
terminology, applying to structure species which are frequently met and
deserve names of their own, even if they turn out to be too restrictive
kind of objects for the use I am having in mind - so why not keep the ter-%
minology already in use, especially for two-groupoids, which is\xout{\unsure{thatxhad}}
pretty well suited after all; b) the structure species I have in mind
is not really a very well determined one, it depends as we saw on choices,
without any one choice looking convincingly better than the others - so
it would be a mess to give an unqualified name to such structure, depen-%
ding on the choice of a certain category $\underline{C}_{\infty} = \underline{C}$ . I have been thinking
of the terminology n-\uline{stack} and $\infty$-\uline{stack} (stack = ``champ'' in French),
a name introduced in Giraud's book (he was restricting to champs =
1-champs ), which over a topos reduced to a one-point space reduces in
his case to the usual notion of a category i.e. 1-category. Here of cour-%
se we are thinking of ``stacks of groupoids'' \unsure{rather|father} than arbitrary stacks,
which I would like to call (for arbitrary order n$\in\underline{N}$ or n=$\infty$) n-Gr-stacks
- suggesting evident ties with the notion of Gr-categories, we should
say Gr-1-categories, of Mme Hoang Xuan Sinh. One advantage of the \deleted{\unsure{namd}}\inserted{name}
``stack'' is that the use it had so far spantaneously suggests the extension
of these notions to the corresponding notions over an arbitrary topos,
which of course is what I am after ultimately. Of course, when an ambi-%
guity is possible, we should speak of \deleted{\illegible{1}} n-$\underline{C}$-stacks - the reference to
$\underline{C}$ should make superfluous the ``Gr'' specification. Thus n-$\underline{C}$-stacks are
essentially the same as n-\unsure{$\underline{C}$}-stacks over the final topos i.e. over a one
point space. When both $\underline{C}$ and ``Gr'' are understood in a given context, we
will use the \xout{\unsure{wordx"stack"}} terminology n-stack simply, or even``stack''
when n is fixed throughout. Thus it will occur that in certain contexts
``stack'' will just mean a usual groupoid, \xout{\unsure{in n=2 groupoid}} in others it
will mean just a category, \xout{\unsure{a n=2-category}} but when n=2 it will not mean
a usual 2-groupoid, but something more general, defined in terms of $\underline{C}$ .

The categories $\underline{C} = \underline{C}_{\infty}$ described before merit a name too - I would
like to call them ``\uline{coherators}'' (``\uline{coh\'ereurs}'' in French), \xout{\illegible{2}}
\over{t}{T}his name is meant to suggest that $\underline{C}$ embodies a hierarchy of coherence
relations, more accurately of coherence``homotopies''. When dealing with
stacks, the term \uline{i-homotopies} (rather than i-objects or i-arrows) for
the elements of the i$^{\mathrm{th}}$ component F$_{i}$ of a stack seems to me the most
suggestive - they will of course be denoted graphically by arrows,

\mspage{\hand{34}}
